% ===========================================================================
%  离散数学与结构（I） Gödel 不完备定理：课堂板书版
%  编码与语法算术化 · 对角线引理 · Gödel/Rosser
%  Compile:  xelatex main   (twice)   or:  latexmk -xelatex main
% ===========================================================================
\documentclass[cjk]{loom}

\newcommand{\goedel}[1]{\ulcorner #1\urcorner}      % Gödel 编码（元语言中的数）
\newcommand{\num}[1]{\overline{#1}}                  % 对象语言中的数词
\newcommand{\N}{\mathbb{N}}
\newcommand{\PA}{\mathsf{PA}}
\newcommand{\Thm}{\mathrm{Thm}}
\newcommand{\Proof}{\mathrm{Proof}}
\newcommand{\Der}{\mathrm{Der}}
\newcommand{\Con}{\mathrm{Con}}
\newcommand{\Term}{\mathrm{Term}}
\newcommand{\Fml}{\mathrm{Fml}}
\newcommand{\Sent}{\mathrm{Sent}}
\newcommand{\Subst}{\mathrm{Subst}}
\newcommand{\Num}{\mathrm{Num}}

\runningthread{离散数学与结构 · Gödel 不完备定理}

\begin{document}

\loomcover
  {Gödel 不完备定理}
  {课堂板书版证明：编码、对角线与 Gödel/Rosser}
  {}
  {2026 年 9 月}

\begin{strand}
本笔记将课堂板书路线整理为一条可检查的形式链：先将语法算术化，使 \(\Thm(y)\) 成为
\(\Sigma_1\) 谓词；再用代入函数证明对角线引理；取 \(\varphi(x)\equiv\neg\Thm(x)\)
得到 Gödel 句，并标明一致性、1-一致性与 Rosser 修改各自承担的作用。
\end{strand}

\clearpage
{\small\tableofcontents}

\clearpage
\section{编码与语法算术化}\label{sec:coding}

\begin{definition}[三层记号]\label{def:three-layers}
\(\goedel{\sigma}\) 是公式或字符串 \(\sigma\) 的 Gödel 编码，属于元语言中的自然数；
\(\num{n}\) 是对象语言中的数词 \(S^{n}0\)。公式内部写 \(\Thm(\goedel{\sigma})\) 时，严格意思是
\(\Thm(\num{\goedel{\sigma}})\)。
\end{definition}

\begin{definition}[常用算术缩写]\label{def:arith-abbrev}
正文直接使用以下正确定义：
\[
x<y:\equiv \exists z\,(y=x+Sz),\qquad
x\mid y:\equiv \exists z\,(y=x\cdot z),
\]
\[
z=x\bmod y:\equiv \exists w\,(x=wy+z\wedge z<y),
\]
\[
\mathrm{Prime}(x):\equiv x>1\wedge \forall y<x\,(y>1\to \neg y\mid x).
\]
这里的全称量词带有上界；有界量化使相应谓词保持为 \(\Delta_0\)，因而“检查到某一位、某一步、某个候选除数”仍是原始递归可判定的。
\end{definition}

\begin{remark}[有限序列必须编码；板书的 \(10^{i}\) 不是系统内免费符号]\label{rem:finite-sequence}
一阶 \(\PA\) 不能直接断称“存在一个有限序列”。该语言只有自然数变量，因此
“存在有限序列 \(a_0,\dots,a_k\)”必须改写为“存在自然数 \(z\)，它编码一个有限序列”，并且
“取第 \(i\) 项”必须可判定。课堂板书的十进制取位
\[
D(N,i)=\left\lfloor \frac{N}{10^{i}}\right\rfloor \bmod 10
\]
在此只能理解为元语言启发：\(\PA\) 的语言没有指数函数符号，直接写 \(10^{i}\) 与直接写 \(2^{i}\) 是同一问题。
若坚持把它内部化，必须先定义 \(\mathrm{Exp10}(i,n)\leftrightarrow n=10^{i}\)；元语言递归
\(10^{0}=1\)、\(10^{i+1}=10\cdot10^{i}\) 仍不够，系统内还要用编码后的有限序列证明
\(\PA\vdash\forall i\,\exists!n\,\mathrm{Exp10}(i,n)\)。因此下文不依赖 \(D(N,i)\)，而抽象假设一个有效序列编码/解码；
它可用 Gödel \(\beta\) 函数实现，\(\beta(c,d,j)=c\bmod(1+(j+1)d)\) 只用乘、加与取余，避免可变幂。
\end{remark}

\begin{definition}[语法谓词与证明谓词]\label{def:syntax-predicates}
假设有效 Gödel 编码已经给定，并且以下函数或谓词原始递归，因而在 Q/\(\PA\) 中可表示：
\(\Term(n)\)、\(\Fml(n)\)、\(\Sent(n)\)、\(\Subst(a,v,t)\)、\(\Num(n)=\goedel{\num n}\)。
证明谓词为
\[
\Proof_{\PA}(p,y):
\]
\(p\) 编码一个有限公式序列，最后一项是 \(y\)，每一项或是逻辑公理/\(\PA\) 公理，或由前面项经推理规则得到。
课堂板书将推导关系记作 \(\Der(x,y,d)\)：可把 \(x\) 理解为编码所用前提或公理序列，把 \(d\) 理解为编码推导；对 \(\PA\) 也可把前提检查折进 \(d\)。
\end{definition}

\begin{definition}[定理谓词]\label{def:thm}
\[
\Thm(y)\equiv \exists p\,\Proof_{\PA}(p,y).
\]
它是 \(\Sigma_1\) 谓词：若 \(y\) 是定理编码，则存在一个有限证明作为可机械检查的见证。
这并不说明定理集可判定，只说明它可半判定。
\end{definition}

\begin{lemma}[\(\Sigma_1\) 完备性的一次使用]\label{lem:d1}
若 \(\sigma\) 有一个标准证明编码 \(p\)，则
\[
\PA\vdash \Proof_{\PA}(\num p,\num{\goedel{\sigma}}),
\qquad
\PA\vdash \Thm(\num{\goedel{\sigma}}).
\]
特别地，
\[
\PA\vdash\sigma \implies \PA\vdash \Thm(\num{\goedel{\sigma}}). \tag{D1}
\]
\end{lemma}

\begin{proof}
\(\Proof_{\PA}\) 是原始递归谓词；对固定的标准 \(p\) 与 \(\goedel{\sigma}\)，命题
\(\Proof_{\PA}(\num p,\num{\goedel{\sigma}})\) 是真 \(\Delta_0\) 句。Q/\(\PA\) 证明所有真
\(\Delta_0\) 句，于是存在量词引入给出 \(\Thm(\num{\goedel{\sigma}})\)。
\end{proof}

\section{对角线引理：把代入做成自指}\label{sec:diagonal}

\begin{definition}[自代入函数]\label{def:self-substitution}
按课堂板书，先将代入写成三元函数 \(\Subst(a,\goedel{x},t)\)，再定义
\[
S(a,b)=\Subst(a,\goedel{x},\Num(b)),\qquad s(u)=S(u,u).
\]
若 \(u=\goedel{\theta(x)}\)，则 \(s(u)=\goedel{\theta(\num u)}\)。函数 \(s\) 原始递归，故 \(\PA\) 可表示并逐个数词计算它。
\end{definition}

\begin{theorem}[对角线引理]\label{thm:diagonal}
对每个只含自由变量 \(x\) 的公式 \(\varphi(x)\)，存在句子 \(\psi\) 使
\[
\PA\vdash \psi\leftrightarrow \varphi(\num{\goedel{\psi}}).
\]
\end{theorem}

\begin{proof}
令
\[
\theta(x)\equiv \varphi(s(x)),\qquad
m=\goedel{\theta(x)},\qquad
\psi\equiv\theta(\num m).
\]
由 \(s\) 的定义，
\[
s(m)=\goedel{\theta(\num m)}=\goedel{\psi}.
\]
因 \(s\) 可表示，\(\PA\) 证明 \(s(\num m)=\num{\goedel{\psi}}\)。于是
\[
\PA\vdash \psi\leftrightarrow \varphi(\num{\goedel{\psi}}).
\]
\end{proof}

\begin{remark}[与课堂板书写法的对应]\label{rem:blackboard-diagonal}
课堂板书写的是
\[
\theta(x)\equiv\varphi(S(x,x)),\quad
m\equiv\goedel{\theta(x)},\quad
\psi\equiv\theta(\num m)
=\varphi(S(\goedel{\theta(x)},\goedel{\theta(x)})).
\]
关键不在代数技巧，而在层次区分：\(\goedel{\theta(x)}\) 是元语言中的数，\(\num m\) 是对象语言中的数词，\(\theta(\num m)\) 又是公式。代入函数作用在编码上，而不是作用在直观文本上。
\end{remark}

\section{Gödel 句与一致性假设}\label{sec:godel}

取
\[
\varphi(x)\equiv\neg\Thm(x).
\]
由定理~\ref{thm:diagonal} 得句子 \(G\) 使
\[
\PA\vdash G\leftrightarrow \neg\Thm(\num{\goedel G}). \tag{FP}
\]
读法：\(G\) 断定“我不是 \(\PA\) 的定理”。

\begin{proposition}[一致性足以排除 \(PA\vdash G\)]\label{prop:not-prove-G}
若 \(\PA\) 一致，则 \(\PA\nvdash G\)。
\end{proposition}

\begin{proof}
若 \(\PA\vdash G\)，由 (D1) 得 \(\PA\vdash\Thm(\num{\goedel G})\)。但 (FP) 给出
\(\PA\vdash\neg\Thm(\num{\goedel G})\)。两者矛盾。
\end{proof}

\begin{proposition}[排除 \(PA\vdash\neg G\) 需要更强的一致性]\label{prop:not-prove-neg-G}
若 \(\PA\) 是 \(\omega\)-一致的，则 \(\PA\nvdash\neg G\)。对本步更弱且够用的情况是 1-一致性：\(\PA\) 证明出的 \(\Sigma_1\) 句在 \(\N\) 中为真。
\end{proposition}

\begin{proof}
若 \(\PA\vdash\neg G\)，由 (FP) 得
\[
\PA\vdash\Thm(\num{\goedel G}),
\]
即
\[
\PA\vdash\exists p\,\Proof_{\PA}(p,\num{\goedel G}). \tag{$\ast$}
\]
一致性本身不能保证 \((\ast)\) 的见证 \(p\) 是标准自然数；它可能只是非标准模型中的见证，在外部并不对应真实有限推导。
\(\omega\)-一致性，或对本步够用的 1-一致性，排除这种非标准见证：此时 \((\ast)\) 给出标准 \(p\)
使 \(\Proof_{\PA}(p,\goedel G)\) 为真，即 \(\PA\vdash G\)，与命题~\ref{prop:not-prove-G} 矛盾。
\end{proof}

\begin{remark}[若课程命题只说“一致性推出不完备”]\label{rem:rosser}
严格地说，应改用 Rosser 句。其核心是把可证性改成“有一个证明我，且没有不晚于它的证明否定我”。
这样普通一致性即可同时排除 \(R\) 与 \(\neg R\)。课堂板书采用的 \(G\) 是 Gödel 原版，因此这里显式记录
\(\omega\)-一致性或 1-一致性这一额外假设。
\end{remark}


\end{document}
